\section{The front end}
\label{sec:frontend}

\subsection{Views and routes}
\label{sec:routes}

Every view and every state has its own link, as in the prototype, so the back button works, a
usability task or a screenshot can start anywhere, and a filter can be shared by copying the address
bar. Filters live in the query string.

{\small
\begin{longtable}{@{}L{0.7cm}L{3.0cm}L{5.7cm}L{2.4cm}L{1.7cm}@{}}
\toprule
\textbf{View} & \textbf{Name} & \textbf{Route} & \textbf{Endpoints} & \textbf{Tier} \\
\midrule
\endhead
V1 & Address selection & \texttt{\#/} \ (\texttt{?q=} prefills the box) & E3, E4 & \must{} (A1) \\
V2 & Address overview & \texttt{\#/address/\{address\}} & E5, E6 & \must{} (A1) \\
V2 & \ldots its transactions tab & \texttt{\#/address/\{address\}/\allowbreak transactions} & E6, E10 & \must \\
V3 & Day and transaction & \texttt{\#/address/\{address\}/\allowbreak day/\{date\}} and \texttt{\ldots/tx/\{txid\}} & E6, E7, E12 & \must{} (A1) \\
V4 & What-if lab & \texttt{\#/address/\{address\}/what-if} & E8, E9 & \must \\
V5 & Models and limits & \texttt{\#/models} & E11 & \must \\
V6 & Watchlist & \texttt{\#/watchlist} & E13 to E15 & \should \\
V7 & About and help & \texttt{\#/about} & E2 & \must \\
V8 & Live lookup & a dialog over V1 & E16 to E19 & \should \\
\bottomrule
\end{longtable}
}

\textbf{The shell around every view.} A header with the product name (a link to V1), the
navigation (Addresses, Models, Watchlist, About), the snapshot date and source, a service-status dot
with its words, and a theme switch. A footer with the limitation line. A skip link, a region for
toasts, and an error boundary. After each navigation the page scrolls to the top and focus moves to
the new heading.

\subsection{V1: address selection}

\begin{itemize}
  \item \textbf{Lookup form.} One field, labelled ``Bitcoin address'', with the hint ``Addresses
    start with 1, 3 or bc1''. The checker runs as the user types and on submit, before any request.
    A valid address goes to E4.
  \item \textbf{Outcomes.} Supported: open V2. Invalid: the specific message under the field, the
    field keeps the text and the focus. Unsupported: a panel that explains the coverage, three
    example addresses as links and ``Clear and try another address''; with S-09, ``Fetch this
    address's public history''.
  \item \textbf{Address table} from E3: address, type, active period, transactions, scored,
    unusual (or ``Short history, not scored''). New: sorting on every column, a search box, and
    filter chips (All, With unusual transactions, Short history, Watched).
  \item \textbf{States.} Skeleton rows while loading; an error panel with ``Try again''; ``No
    address matches; clear the filter'' when a filter empties the table.
  \item \textbf{With S-04:} a star on each row and up to five ``Recently opened'' chips.
\end{itemize}

\subsection{V2: address overview}

\begin{itemize}
  \item \textbf{Heading.} The full address with a copy button, its type, the source and snapshot
    date, an Export menu (CSV, JSON), and with S-01 the method selector.
  \item \textbf{Four summary tiles}: transactions, active period, unusual transactions and days,
    and ``Scored 367 of 377''.
  \item \textbf{Tabs}: Overview, Transactions, What-if and, with S-06, Replay. Each tab is a route.
  \item \textbf{Overview tab.} Chart G1 and the list ``Days with unusual transactions'', which
    shows eight days and a ``Show all'' control (the fix for Assignment 1's inspection issue 2).
  \item \textbf{Transactions tab.} The filter bar, chart G2 and the transaction table
    (Section~\ref{sec:filters}).
  \item \textbf{Short history.} The banner from the prototype with its two actions, ``Browse its
    transactions anyway'' and ``Choose an address with more history''; the tiles and G1 still show
    the counts.
\end{itemize}

\subsection{The filter bar and the transaction table (M-08)}
\label{sec:filters}

\begin{itemize}
  \item \textbf{Controls.} Status chips (Unusual, Not unusual, Not scored), direction chips
    (Received, Sent), a date range (two date fields limited to the address's active period), a
    sort menu (newest, oldest, largest, highest score) and ``Reset filters''.
  \item \textbf{Behaviour.} Each change updates the query string, sends E6 and redraws chart G2
    and the table together. The count line says ``Showing 12 of 377 transactions''. Zooming G2 on
    the time axis fills the date range; the two stay in step.
  \item \textbf{Table.} Time (UTC), direction, amount in BTC to eight places, score, status badge
    and the main reason; 50 rows a page; a row opens V3. On a phone each row becomes a card.
  \item \textbf{Validation.} A ``from'' date after the ``to'' date marks both fields and sends
    nothing. The server applies the same rule.
\end{itemize}

\subsection{V3: day and transaction}

\begin{itemize}
  \item \textbf{Layout.} The day's transactions on the left and the detail panel on the right. On
    a phone the panel opens as a full-height sheet over the list, with a Close button and focus held
    inside it.
  \item \textbf{Facts (FR-03).} Transaction ID with a copy button, confirmed time in UTC, block,
    direction, the amount and how it is made up from received and sent, and a ``View on
    Blockstream'' link that opens in a new tab.
  \item \textbf{Why this transaction stands out (FR-04, AC-03).} The explanation in plain words,
    the comparison basis (``10 earlier transactions in the same direction, 16 May to 9 July 2024''),
    and the limitation line. With S-05, the Assistant control of Section~\ref{sec:assistant}.
  \item \textbf{Chart G3} and the table that repeats every number. The z-score appears in the
    table and the tooltip, not in the sentence (inspection issue 1).
  \item \textbf{What each model says (M-13).} A table of the nine methods: family, score,
    threshold and verdict (``Flag'', a dash, or ``not available''), with the deciding method
    marked. With S-02, chart G4 above it.
  \item \textbf{Moving on.} Previous and next day, previous and next unusual day, and ``Try a
    what-if from here'', which opens V4 with this transaction's direction.
  \item \textbf{States.} Nothing selected: ``Choose a transaction to see its details''. Not
    scored: the facts, then the insufficient-history note with its two next steps.
\end{itemize}

\subsection{V4: the what-if lab (M-06)}
\label{sec:viewwhatif}

This is the brief's ``custom user input form with comprehensive validation and dynamic feedback''
and its ``real-time prediction''.

\begin{longtable}{@{}L{2.2cm}L{5.6cm}L{6.2cm}@{}}
\toprule
\textbf{Field} & \textbf{Control} & \textbf{Feedback while typing} \\
\midrule
\endhead
Amount & A text field in decimal mode with a BTC or satoshi switch; a slider on a log scale over the curve's range; chips for ``Usual'', ``10 × usual'' and ``100 × usual'' & The rule that fails, under the field; else a line such as ``0.05 BTC is about 95 times this address's usual outgoing amount'' \\
Direction & A two-way switch: Received, Sent & The usual amount and the comparison count change with it \\
Time (UTC) & A date field and a time field; chips for ``The day after its last transaction'', ``After 30 quiet days'' and ``After 90 quiet days'' & The gap since the previous transaction, in days and hours; a message if the time is before the first transaction or too far ahead \\
\bottomrule
\end{longtable}

\begin{itemize}
  \item \textbf{No submit button is needed.} When all three fields are valid, the request goes
    250 ms after the last change and any older request is cancelled. A ``Score it'' button is still
    there for keyboard users and for anyone who expects one.
  \item \textbf{While it updates.} The previous result stays in place, dimmed, with an ``Updating''
    mark; nothing jumps.
  \item \textbf{The result.} A verdict card (``Flagged by Robust z baseline'' or ``Not flagged'',
    the score against the flag line, the explanation, and the limitation line when flagged); chart
    G3 for the hypothetical transaction; the table of eight models, with TxGPT's row marked ``not
    available for a hypothetical transaction''; chart G5 with the current amount marked.
  \item \textbf{Under the curve}, one sentence read from it. For the running example it is: ``On
    this address, a sent amount would be flagged from about 0.023 BTC upwards, and below about
    0.00001 BTC.''
  \item \textbf{A note at the top:} ``A what-if is hypothetical. Nothing is sent to the Bitcoin
    network and nothing is stored.''
  \item \textbf{Errors.} 422: ``At that time this address had 6 earlier transactions; scoring
    needs 10 over 14 days'', with a button that sets the first date that works. A network failure
    keeps every value in the form.
  \item \textbf{With S-08:} ``Keep this scenario'' adds a row to a table of up to ten scenarios
    (amount, direction, time, verdict, score, how many models flag it), which can be exported.
\end{itemize}

\subsection{V5 to V8}
\label{sec:viewwatch}

\begin{itemize}
  \item \textbf{V5, models and limits.} A card per method from E11: name, family, one sentence on
    how it scores, its threshold and whether it can score new input. With S-03, chart G6. Then
    ``What a flag means, and what it does not'': flags compare an address with its own past; real
    transactions have no labels, so no accuracy is claimed for them; recall is measured on planted
    test transactions; the snapshot covers 28 addresses on one date.
  \item \textbf{V6, watchlist (S-04).} Entries with label, note and counts. The star on V1 and V2
    opens a dialog with a label field and a note field, each with a character count. Removing an
    entry takes effect at once and shows a toast with Undo.
  \item \textbf{V7, about and help.} What the Explorer is for and not for; how to read a flag; the
    four features in plain words; the data source, snapshot and terms; keyboard help; the version.
  \item \textbf{V8, live lookup (S-09).} A dialog with a progress bar, the page and transaction
    counts and Cancel. Rate-limited: the wait, counting down. Upstream failure: ``Try again'' and
    ``Pick a snapshot address''.
\end{itemize}

\subsection{State and data}

\begin{itemize}
  \item \textbf{Server data} comes through \code{useRequest}, keyed by the request's address. The
    hook gains cancellation, an optional refresh interval, and ``keep the previous data while
    loading''. Snapshot answers are kept for the session.
  \item \textbf{The address bar} is the source of truth for the view, the selected transaction
    and the filters.
  \item \textbf{Preferences} (theme, units, recent addresses, the last method and filters) are in
    local storage under one versioned key.
  \item \textbf{The watchlist} is on the server. \textbf{Scenarios} are in session storage.
  \item \textbf{Settings} come from E2 once at start. No store library is needed.
\end{itemize}

\subsection{Components}
\label{sec:components}

\begin{center}
\small
\begin{tabularx}{\textwidth}{@{}L{3.3cm}Y@{}}
\toprule
\textbf{From the prototype} & Badges, breadcrumbs, copy button, error state, state panel, page heading, icons \\
\textbf{New, general} & Tabs; modal dialog (focus held inside, Escape closes, focus returns); toast with Undo; tooltip; skeleton; progress bar; data table with sorting and paging; empty state; error boundary; theme switch; service-status indicator \\
\textbf{New, forms} & Field (label, hint, error, described-by links); number field with a unit switch; log slider; two-way switch; date and time fields; character counter; filter bar with chips; method menu; export menu \\
\bottomrule
\end{tabularx}
\end{center}

The rubric's examples of advanced components are modals and dynamic forms. The build has three
dialogs (watchlist note, live lookup, the phone's detail sheet), one dynamic form (V4), tabs, toasts
with Undo and skeleton loading.

\subsection{Layouts and touch}
\label{sec:responsive}

\begin{center}
\small
\begin{tabularx}{\textwidth}{@{}L{2.3cm}Y Y Y@{}}
\toprule
 & \textbf{Phone, 390 px} & \textbf{Tablet, 820 px} & \textbf{Desktop, 1280 px} \\
\midrule
Navigation & A menu button & A row & A row \\
V1 table & Cards & Table, fewer columns & Full table \\
V2 tiles & Two by two & Four in a row & Four in a row \\
V3 & List, then a sheet & List above the panel & List beside the panel \\
V4 & Form, then result & Form beside the verdict & Form, verdict and curve in view together \\
Charts & Full width; fewer ticks; tools in a menu & Full width & Up to two side by side \\
\bottomrule
\end{tabularx}
\end{center}

No page scrolls sideways; the screenshot script fails if one does. Touch targets are at least
44 px. On a touch screen a tap shows a chart's tooltip and a second tap opens the item.

\subsection{Accessibility (NFR-03)}
\label{sec:a11y}

\begin{itemize}
  \item Every action is a link or a button, reachable by Tab in reading order, with a visible
    focus ring.
  \item Chart G1 keeps the prototype's keyboard control: arrow keys move between days, Enter
    opens one. Every chart has a text summary and a ``View as table'' switch.
  \item A flag is never colour alone: orange, a triangle and the word ``Unusual''.
  \item Text contrast is at least 4.5 to 1 in both themes; the chart colours are checked for
    colour-blind readers.
  \item Form errors are tied to their field, announced when they appear, and never clear what the
    user typed.
  \item Toasts and loading messages use a polite live region. Motion stops when the system asks
    for reduced motion. The layout holds at 200\% zoom.
\end{itemize}

\subsection{Experience details (Usability 2)}

\begin{itemize}
  \item \textbf{Feedback.} Something visible within 100 ms of every action: a pressed state, a
    skeleton, a progress line. Success is confirmed (``Copied'', ``Exported 377 rows'', ``Added to
    the watchlist''). Every error says what happened and offers a next step.
  \item \textbf{Personal.} The theme, the units, the last method and filters and the recent
    addresses are remembered; the watchlist carries the user's own labels and notes.
  \item \textbf{Transitions.} Panels, dialogs and toasts ease in and out in 150 to 200 ms; chart
    updates animate between states and keep the zoom.
  \item \textbf{Words.} One vocabulary everywhere (Section~\ref{sec:wording}).
\end{itemize}

\subsection{Wording}
\label{sec:wording}

\begin{center}
\small
\begin{tabularx}{\textwidth}{@{}L{4.2cm}Y@{}}
\toprule
\textbf{Say} & unusual; flag, flagged; worth a look; stands out; compared with this address's own history; not scored \\
\textbf{Never say} & fraud, criminal, illicit, guilty, suspicious, for the tool or for any address (\code{bannedWords} in the settings); also no ``risk score'', ``alert'', ``threat'' or ``safe'' for an address \\
\textbf{Always show} & ``Unusual activity requires investigation; it does not indicate wrongdoing.'' on every flagged item, in the footer, in every export \\
\bottomrule
\end{tabularx}
\end{center}

A test renders every view and state and fails on a banned word; the back end's test does the same
for every string in every response.

\subsection{The usability walkthrough and the accessibility inspection (M-15)}
\label{sec:usability}

\begin{itemize}
  \item \textbf{Who and when.} Two or more people who did not build the application; Wednesday 25
    and Thursday 26 November, on the MS2 build. Minh Khang Nguyen runs the sessions; another member
    observes with Assignment 1's observer sheet.
  \item \textbf{Tasks.} Assignment 1's four, unchanged: (1) open the activity of a given address;
    (2) find a day with unusual activity and open it; (3) find the transaction that stands out and
    say what the marker means; (4) look up an address outside the snapshot and decide what to do.
    A fifth is added: (5) find out how large a payment would have to be to stand out on this
    address.
  \item \textbf{Success.} Task 3 passes when the participant describes the flag as ``worth a
    closer look'' and not as wrongdoing (NFR-01, condition 4). Each task records completion, time,
    errors and the participant's words.
  \item \textbf{After.} Issues are rated on Nielsen's 0 to 4 scale; every 3 and 4 is fixed before
    MS4 and retested.
  \item \textbf{Accessibility inspection.} J1 to J4 by keyboard alone; the contrast check; a pass
    with a screen reader; reduced motion; 200\% zoom. Findings go in the same log.
\end{itemize}

\begin{done}
\fp{docs/usability/} holds the filled observer sheets and the issue log; the report's appendix has
the results table; no issue rated 3 or 4 is open.
\end{done}
